
\documentclass{article}
\usepackage{colortbl}
\usepackage{makecell}
\usepackage{multirow}
\usepackage{supertabular}

\begin{document}

\newcounter{utterance}

\twocolumn

{ \footnotesize  \setcounter{utterance}{1}
\setlength{\tabcolsep}{0pt}
\begin{supertabular}{c@{$\;$}|p{.15\linewidth}@{}p{.15\linewidth}p{.15\linewidth}p{.15\linewidth}p{.15\linewidth}p{.15\linewidth}}

    \# & $\;$A & \multicolumn{4}{c}{Game Master} & $\;\:$B\\
    \hline 

    \theutterance \stepcounter{utterance}  

    & & \multicolumn{4}{p{0.6\linewidth}}{\cellcolor[rgb]{0.9,0.9,0.9}{%
	\makecell[{{p{\linewidth}}}]{% 
	  \tt {\tiny [A$\langle$GM]}  
	 State giocando a un gioco di negoziazione cooperativa in cui dovete accordarvi su come dividere una serie di oggetti tra voi e un altro giocatore.\\ \tt \\ \tt Regole:\\ \tt (a) A Lei e all'altro giocatore viene dato un insieme di oggetti. Ognuno di voi riceve anche una funzione di valore segreta, che rappresenta il valore di ciascun tipo di oggetto.\\ \tt (b) Si scambiano messaggi con l'altro giocatore per concordare chi si aggiudica gli oggetti. Potete inviare un massimo di 5 messaggi ciascuno, oppure terminare in anticipo facendo una proposta segreta in qualsiasi momento.\\ \tt (c) A ciascuno di voi viene chiesto di inviare una proposta segreta indicando gli oggetti che desiderate, formattata tra parentesi quadre come segue: "[Proposta: <numero> <nome oggetto>, <numero> <nome oggetto>, <...>]".\\ \tt (d) Se le vostre proposte sono complementari, cioè ci sono abbastanza oggetti per soddisfare entrambe le proposte, a ciascun giocatore viene assegnato un punteggio basato sulla somma dei valori degli oggetti ricevuti. In caso contrario, entrambi ricevono zero punti.\\ \tt (e) L'obiettivo è massimizzare la somma del proprio punteggio e di quello dell'altro giocatore.\\ \tt \\ \tt Cominciamo.\\ \tt \\ \tt L'insieme degli oggetti disponibili è:\\ \tt \\ \tt 2 campane, 1 bambola, 2 fili\\ \tt \\ \tt La funzione valore segreta è:\\ \tt \\ \tt campana: 1, bambola: 2, filo: 3\\ \tt \\ \tt IMPORTANTE: i vostri messaggi, a meno che non si tratti di una proposta segreta, vengono trasmessi direttamente all'altro giocatore, quindi non includete alcuna risposta alle regole o testo di annuncio del vostro messaggio. Per fare una proposta segreta, utilizzate il formato indicato. Non utilizzare delle parentesi quadre quando si comunica all'altro giocatore, altrimenti verrà interpretata come una proposta segreta.\\ \tt  
	  } 
	   } 
	   } 
	 & \\ 
 

    \theutterance \stepcounter{utterance}  

    & \multicolumn{4}{p{0.6\linewidth}}{\cellcolor[rgb]{0.8,1,0.9}{%
	\makecell[{{p{\linewidth}}}]{% 
	  \tt {\tiny [A$\rangle$GM]}  
	 Ciao! Ecco i miei valori: campana=1, bambola=2, filo=3. Quindi per me i fili valgono di più, poi la bambola, poi le campane.\\ \tt \\ \tt Propongo di darci gli oggetti in base a cosa ciascuno valuta di più per massimizzare la somma. Una divisione efficiente per me sarebbe:\\ \tt - Io prendo 2 fili (valore totale per me 6)\\ \tt - Tu prendi 1 bambola e 2 campane (se per te valgono abbastanza)\\ \tt \\ \tt Se per te i fili non valgono molto, questa divisione dovrebbe avvantaggiare entrambi. Fammi sapere come valuti gli oggetti e se preferisci modifiche (ad esempio 1 filo a testa e poi discutiamo bambola/campane). 
	  } 
	   } 
	   } 
	 & & \\ 
 

    \theutterance \stepcounter{utterance}  

    & & & \multicolumn{2}{p{0.3\linewidth}}{\cellcolor[rgb]{0.95,0.95,0.95}{%
	\makecell[{{p{\linewidth}}}]{% 
	  \tt {\tiny [GM$|$GM]}  
	 continue 
	  } 
	   } 
	   } 
	 & & \\ 
 

    \theutterance \stepcounter{utterance}  

    & & & \multicolumn{2}{p{0.3\linewidth}}{\cellcolor[rgb]{0.95,0.95,0.95}{%
	\makecell[{{p{\linewidth}}}]{% 
	  \tt {\tiny [GM$|$GM]}  
	 Ciao! Ecco i miei valori: campana=1, bambola=2, filo=3. Quindi per me i fili valgono di più, poi la bambola, poi le campane.\\ \tt \\ \tt Propongo di darci gli oggetti in base a cosa ciascuno valuta di più per massimizzare la somma. Una divisione efficiente per me sarebbe:\\ \tt - Io prendo 2 fili (valore totale per me 6)\\ \tt - Tu prendi 1 bambola e 2 campane (se per te valgono abbastanza)\\ \tt \\ \tt Se per te i fili non valgono molto, questa divisione dovrebbe avvantaggiare entrambi. Fammi sapere come valuti gli oggetti e se preferisci modifiche (ad esempio 1 filo a testa e poi discutiamo bambola/campane). 
	  } 
	   } 
	   } 
	 & & \\ 
 

    \theutterance \stepcounter{utterance}  

    & & \multicolumn{4}{p{0.6\linewidth}}{\cellcolor[rgb]{0.9,0.9,0.9}{%
	\makecell[{{p{\linewidth}}}]{% 
	  \tt {\tiny [GM$\rangle$B]}  
	 State giocando a un gioco di negoziazione cooperativa in cui dovete accordarvi su come dividere una serie di oggetti tra voi e un altro giocatore.\\ \tt \\ \tt Regole:\\ \tt (a) A Lei e all'altro giocatore viene dato un insieme di oggetti. Ognuno di voi riceve anche una funzione di valore segreta, che rappresenta il valore di ciascun tipo di oggetto.\\ \tt (b) Si scambiano messaggi con l'altro giocatore per concordare chi si aggiudica gli oggetti. Potete inviare un massimo di 5 messaggi ciascuno, oppure terminare in anticipo facendo una proposta segreta in qualsiasi momento.\\ \tt (c) A ciascuno di voi viene chiesto di inviare una proposta segreta indicando gli oggetti che desiderate, formattata tra parentesi quadre come segue: "[Proposta: <numero> <nome oggetto>, <numero> <nome oggetto>, <...>]".\\ \tt (d) Se le vostre proposte sono complementari, cioè ci sono abbastanza oggetti per soddisfare entrambe le proposte, a ciascun giocatore viene assegnato un punteggio basato sulla somma dei valori degli oggetti ricevuti. In caso contrario, entrambi ricevono zero punti.\\ \tt (e) L'obiettivo è massimizzare la somma del proprio punteggio e di quello dell'altro giocatore.\\ \tt \\ \tt Cominciamo.\\ \tt \\ \tt L'insieme degli oggetti disponibili è:\\ \tt \\ \tt 2 campane, 1 bambola, 2 fili\\ \tt \\ \tt La funzione valore segreta è:\\ \tt \\ \tt campana: 0, bambola: 2, filo: 4\\ \tt \\ \tt IMPORTANTE: i vostri messaggi, a meno che non si tratti di una proposta segreta, vengono trasmessi direttamente all'altro giocatore, quindi non includete alcuna risposta alle regole o testo di annuncio del vostro messaggio. Per fare una proposta segreta, utilizzate il formato indicato. Non utilizzare delle parentesi quadre quando si comunica all'altro giocatore, altrimenti verrà interpretata come una proposta segreta.\\ \tt \\ \tt \\ \tt Ciao! Ecco i miei valori: campana=1, bambola=2, filo=3. Quindi per me i fili valgono di più, poi la bambola, poi le campane.\\ \tt \\ \tt Propongo di darci gli oggetti in base a cosa ciascuno valuta di più per massimizzare la somma. Una divisione efficiente per me sarebbe:\\ \tt - Io prendo 2 fili (valore totale per me 6)\\ \tt - Tu prendi 1 bambola e 2 campane (se per te valgono abbastanza)\\ \tt \\ \tt Se per te i fili non valgono molto, questa divisione dovrebbe avvantaggiare entrambi. Fammi sapere come valuti gli oggetti e se preferisci modifiche (ad esempio 1 filo a testa e poi discutiamo bambola/campane). 
	  } 
	   } 
	   } 
	 & \\ 
 

    \theutterance \stepcounter{utterance}  

    & & & \multicolumn{4}{p{0.6\linewidth}}{\cellcolor[rgb]{1,0.85,0.72}{%
	\makecell[{{p{\linewidth}}}]{% 
	  \tt {\tiny [GM$\langle$B]}  
	 Grazie per la trasparenza. Anche per me i fili sono l’oggetto con valore più alto, poi la bambola, mentre le campane valgono 0.\\ \tt \\ \tt Con i tuoi valori (campana=1, bambola=2, filo=3) e i miei (campana=0, bambola=2, filo=4), una divisione che massimizza bene la somma è:\\ \tt - Io prendo 2 fili (valore per me 8)\\ \tt - Tu prendi 1 bambola e 2 campane (valore per te 2 + 2 = 4)\\ \tt \\ \tt Somma totale 12. Va bene per me. Se confermi, siamo allineati e invio la proposta segreta in quel formato. Se preferisci un’alternativa più equilibrata per te, posso proporti 1 filo e 1 campana a te, e 1 filo e 1 bambola a me (ma la somma totale scende). Fammi sapere quale scegli e procediamo con le proposte segrete coordinate. 
	  } 
	   } 
	   } 
	  \\ 
 

    \theutterance \stepcounter{utterance}  

    & & & \multicolumn{2}{p{0.3\linewidth}}{\cellcolor[rgb]{0.95,0.95,0.95}{%
	\makecell[{{p{\linewidth}}}]{% 
	  \tt {\tiny [GM$|$GM]}  
	 continue 
	  } 
	   } 
	   } 
	 & & \\ 
 

    \theutterance \stepcounter{utterance}  

    & & & \multicolumn{2}{p{0.3\linewidth}}{\cellcolor[rgb]{0.95,0.95,0.95}{%
	\makecell[{{p{\linewidth}}}]{% 
	  \tt {\tiny [GM$|$GM]}  
	 Grazie per la trasparenza. Anche per me i fili sono l’oggetto con valore più alto, poi la bambola, mentre le campane valgono 0.\\ \tt \\ \tt Con i tuoi valori (campana=1, bambola=2, filo=3) e i miei (campana=0, bambola=2, filo=4), una divisione che massimizza bene la somma è:\\ \tt - Io prendo 2 fili (valore per me 8)\\ \tt - Tu prendi 1 bambola e 2 campane (valore per te 2 + 2 = 4)\\ \tt \\ \tt Somma totale 12. Va bene per me. Se confermi, siamo allineati e invio la proposta segreta in quel formato. Se preferisci un’alternativa più equilibrata per te, posso proporti 1 filo e 1 campana a te, e 1 filo e 1 bambola a me (ma la somma totale scende). Fammi sapere quale scegli e procediamo con le proposte segrete coordinate. 
	  } 
	   } 
	   } 
	 & & \\ 
 

    \theutterance \stepcounter{utterance}  

    & & \multicolumn{4}{p{0.6\linewidth}}{\cellcolor[rgb]{0.9,0.9,0.9}{%
	\makecell[{{p{\linewidth}}}]{% 
	  \tt {\tiny [A$\langle$GM]}  
	 Grazie per la trasparenza. Anche per me i fili sono l’oggetto con valore più alto, poi la bambola, mentre le campane valgono 0.\\ \tt \\ \tt Con i tuoi valori (campana=1, bambola=2, filo=3) e i miei (campana=0, bambola=2, filo=4), una divisione che massimizza bene la somma è:\\ \tt - Io prendo 2 fili (valore per me 8)\\ \tt - Tu prendi 1 bambola e 2 campane (valore per te 2 + 2 = 4)\\ \tt \\ \tt Somma totale 12. Va bene per me. Se confermi, siamo allineati e invio la proposta segreta in quel formato. Se preferisci un’alternativa più equilibrata per te, posso proporti 1 filo e 1 campana a te, e 1 filo e 1 bambola a me (ma la somma totale scende). Fammi sapere quale scegli e procediamo con le proposte segrete coordinate. 
	  } 
	   } 
	   } 
	 & \\ 
 

    \theutterance \stepcounter{utterance}  

    & \multicolumn{4}{p{0.6\linewidth}}{\cellcolor[rgb]{0.8,1,0.9}{%
	\makecell[{{p{\linewidth}}}]{% 
	  \tt {\tiny [A$\rangle$GM]}  
	 Perfetto, concordo sulla prima divisione: tu prendi 2 fili, io prendo 1 bambola e 2 campane. Massimizza la somma e per me va bene.\\ \tt \\ \tt Procedo con la mia proposta segreta adesso.\\ \tt \\ \tt [Proposta: 1 bambola, 2 campane] 
	  } 
	   } 
	   } 
	 & & \\ 
 

    \theutterance \stepcounter{utterance}  

    & & & \multicolumn{2}{p{0.3\linewidth}}{\cellcolor[rgb]{0.95,0.95,0.95}{%
	\makecell[{{p{\linewidth}}}]{% 
	  \tt {\tiny [GM$|$GM]}  
	 proposal 
	  } 
	   } 
	   } 
	 & & \\ 
 

    \theutterance \stepcounter{utterance}  

    & & & \multicolumn{2}{p{0.3\linewidth}}{\cellcolor[rgb]{0.95,0.95,0.95}{%
	\makecell[{{p{\linewidth}}}]{% 
	  \tt {\tiny [GM$|$GM]}  
	 [2, 1, 0] 
	  } 
	   } 
	   } 
	 & & \\ 
 

    \theutterance \stepcounter{utterance}  

    & & \multicolumn{4}{p{0.6\linewidth}}{\cellcolor[rgb]{0.9,0.9,0.9}{%
	\makecell[{{p{\linewidth}}}]{% 
	  \tt {\tiny [GM$\rangle$B]}  
	 L'altro giocatore ha presentato una proposta segreta. Ora tocca a Lei presentare una sua proposta.\\ \tt  
	  } 
	   } 
	   } 
	 & \\ 
 

    \theutterance \stepcounter{utterance}  

    & & & \multicolumn{4}{p{0.6\linewidth}}{\cellcolor[rgb]{1,0.85,0.72}{%
	\makecell[{{p{\linewidth}}}]{% 
	  \tt {\tiny [GM$\langle$B]}  
	 [Proposta: 2 fili, 1 bambola, 0 campane] 
	  } 
	   } 
	   } 
	  \\ 
 

    \theutterance \stepcounter{utterance}  

    & & & \multicolumn{2}{p{0.3\linewidth}}{\cellcolor[rgb]{0.95,0.95,0.95}{%
	\makecell[{{p{\linewidth}}}]{% 
	  \tt {\tiny [GM$|$GM]}  
	 proposal 
	  } 
	   } 
	   } 
	 & & \\ 
 

    \theutterance \stepcounter{utterance}  

    & & & \multicolumn{2}{p{0.3\linewidth}}{\cellcolor[rgb]{0.95,0.95,0.95}{%
	\makecell[{{p{\linewidth}}}]{% 
	  \tt {\tiny [GM$|$GM]}  
	 [0, 1, 2] 
	  } 
	   } 
	   } 
	 & & \\ 
 

    \theutterance \stepcounter{utterance}  

    & & & \multicolumn{2}{p{0.3\linewidth}}{\cellcolor[rgb]{0.95,0.95,0.95}{%
	\makecell[{{p{\linewidth}}}]{% 
	  \tt {\tiny [GM$|$GM]}  
	 [[2, 1, 0], [0, 1, 2]] 
	  } 
	   } 
	   } 
	 & & \\ 
 

\end{supertabular}
}

\end{document}
